\documentclass[11pt,a4paper]{article}
\usepackage[margin=2.5cm]{geometry}
\usepackage{booktabs}
\usepackage{amsmath}
\usepackage{hyperref}

\title{Finite-Volume Voxel-Spacing-Aware GLDM: Physical Rationale and Synthetic Validation}
\author{}
\date{}

\begin{document}
\maketitle

\section*{Purpose}
This note documents a refinement of the voxel-spacing-aware GLDM implementation in the corrected modified PyRadiomics prototype. The previous GLDM prototype used a weighted dependence size and then discretized that value before inserting it into the GLDM matrix. Although useful as an experimental proof of concept, that approach changed the interpretation of the GLDM dependence axis.

The new implementation treats anisotropic voxels as finite physical support elements and evaluates standard GLDM dependence counting on a zero-order-hold physical lattice. This keeps dependence size as an integer count and avoids fractional or rounded dependence sizes.

\section*{Physical Motivation}
An anisotropic voxel should not be interpreted only as a point sample. For example, with spacing
\[
(\Delta_x,\Delta_y,\Delta_z)=(1,1,5)\,\mathrm{mm},
\]
a through-plane voxel represents a thick physical support. A reasonable finite-volume interpretation is that this voxel corresponds to five \(1\times1\times1\) mm subcells with the same native intensity value. This is not interpolation: no new intensity gradients are introduced. It is a zero-order-hold representation of the measured voxel support.

\section*{Mathematical Definition}
Let \(I\) denote the native binned image and let \(s=(\Delta_z,\Delta_y,\Delta_x)\) be the voxel spacing in array axis order. Define
\[
h = \min(\Delta_z,\Delta_y,\Delta_x),
\]
and integer repeat factors
\[
r_k = \mathrm{round}\left(\frac{\Delta_k}{h}\right).
\]
The finite-volume image \(\tilde I\) is produced by repeating each native voxel \(r_z\times r_y\times r_x\) times:
\[
\tilde I = \mathrm{ZOH}(I,s),
\]
where ZOH denotes zero-order hold.

Standard GLDM dependence counting is then performed on \(\tilde I\):
\[
j(\tilde x)=
\sum_{b\in N_{\delta}^{std}}
\mathbf{1}\left(|\tilde I(\tilde x)-\tilde I(\tilde x+b)|\leq\alpha\right),
\]
and the GLDM matrix is updated as
\[
P(\tilde I(\tilde x),j(\tilde x)) \leftarrow
P(\tilde I(\tilde x),j(\tilde x))+1.
\]

Thus, \(j\) remains an integer dependence count and \(P(i,j)\) remains a count matrix. The modification changes the physical support on which the standard GLDM neighborhood is evaluated, not the GLDM dependence axis itself.

\section*{Implementation}
The implementation was applied only to:
\[
\texttt{<REPOSITORY_ROOT>/librerias/radiomics\_modificado\_corregido/gldm.py}.
\]
When \texttt{weightingNorm = voxelSpacing} and spacing is anisotropic, the binned image array and mask are expanded by integer repeat factors derived from spacing ratios, then the standard C GLDM backend is called. When spacing is isotropic, the implementation routes to the standard GLDM path exactly.

\section*{Validation}
The synthetic validation used a deterministic anisotropic texture image with SimpleITK spacing \((1.0,1.0,5.0)\), corresponding to array-order spacing \((5.0,1.0,1.0)\), and an isotropic sanity image with spacing \((1.0,1.0,1.0)\). Extractions were run in isolated Python 3.10 arm64 environments with freshly compiled C extensions.

\begin{table}[h]
\centering
\caption{Global validation checks.}
\begin{tabular}{lrrrr}
\toprule
Comparison & Features & Mean abs. diff. & Max abs. diff. & Changed \\
\midrule
Original anisotropic vs modified default anisotropic & 75 & 0.000000 & 0.000000 & 0 \\
Modified default isotropic vs modified voxelSpacing isotropic & 75 & 0.000000 & 0.000000 & 0 \\
\bottomrule
\end{tabular}
\end{table}

\begin{table}[h]
\centering
\caption{Anisotropic activation by family after finite-volume GLDM refinement.}
\begin{tabular}{lrrrr}
\toprule
Family & Features & Changed & Median abs. diff. & Median rel. diff. \\
\midrule
GLCM & 24 & 24 & 0.095800 & 0.101046 \\
GLRLM & 16 & 16 & 0.094053 & 0.139830 \\
NGTDM & 5 & 5 & 0.022053 & 0.335726 \\
GLDM & 14 & 11 & 0.117819 & 0.491003 \\
GLSZM & 16 & 0 & 0.000000 & 0.000000 \\
\bottomrule
\end{tabular}
\end{table}

\begin{table}[h]
\centering
\caption{Selected GLDM feature changes under anisotropic voxelSpacing mode.}
\begin{tabular}{lrrr}
\toprule
Feature & Default & Finite-volume voxelSpacing & Absolute diff. \\
\midrule
DependenceEntropy & 6.327318 & 6.388388 & 0.061070 \\
DependenceVariance & 16.546574 & 38.996144 & 22.449570 \\
GrayLevelNonUniformity & 1183.017072 & 5915.085359 & 4732.068287 \\
LargeDependenceEmphasis & 163.478009 & 331.453646 & 167.975637 \\
SmallDependenceEmphasis & 0.014030 & 0.007205 & 0.006825 \\
\bottomrule
\end{tabular}
\end{table}

\section*{Interpretation}
The refined GLDM behaves coherently in the synthetic validation. Default extraction remains exactly backward compatible with original PyRadiomics. The isotropic sanity check remains exact, indicating that the finite-volume path does not introduce an artificial change when voxel dimensions are equal. Under anisotropic spacing, 11 of 14 GLDM features changed.

Importantly, GLDM gray-level-only emphasis features remained unchanged in the synthetic comparison, including GrayLevelVariance, HighGrayLevelEmphasis and LowGrayLevelEmphasis. This is physically reassuring: the finite-volume modification affects dependence structure and local spatial support rather than altering the gray-level distribution itself.

\section*{Manuscript Wording}
A concise manuscript description could be:

\begin{quote}
For GLDM, anisotropic voxels were treated as finite physical support elements rather than point samples. Under voxel-spacing-aware extraction, the discretized native image and mask were expanded by zero-order hold onto a physical lattice with step equal to the smallest voxel spacing. Standard GLDM dependence counting was then applied on this finite-volume representation. This preserves the integer dependence-size axis and count-matrix interpretation of GLDM while allowing anisotropic voxel support to influence local dependence structure without intensity interpolation.
\end{quote}

\section*{Limitations}
This remains an extension of standard GLDM rather than an IBSI-defined reference implementation. For non-integer spacing ratios, repeat factors are rounded and a warning is issued. The approach can increase computation and memory use when one axis is much thicker than the others. However, for common anisotropic CT/MRI settings such as \(1\times1\times5\) mm, the expansion is straightforward and remains conceptually faithful to the finite voxel support.

\end{document}
